% 文献翻译流程测试文件
% 基于 GDUT 模板最小化改写，仅验证 skill 流程，不翻译整篇。

\documentclass[12pt,a4paper]{article}
\usepackage[UTF8]{ctex}
\usepackage{ifxetex,ifluatex}
\ifnum\ifxetex1\else0\fi=1
    \usepackage{fontspec}
\else
    \ifnum\ifluatex1\else0\fi=1
        \usepackage{fontspec}
    \else
        \usepackage{times}
    \fi
\fi
\usepackage{geometry}
\usepackage{titlesec}
\usepackage{fancyhdr}
\usepackage{setspace}
\usepackage{graphicx}
\usepackage{amsmath}
\usepackage{amssymb}
\usepackage{booktabs}
\usepackage{array}
\usepackage{multirow}
\usepackage{longtable}
\usepackage{cite}
\usepackage{enumitem}
\usepackage{indentfirst}
\usepackage{float}
\usepackage{caption}
\usepackage{subcaption}
\usepackage{hyperref}
\usepackage{pdfpages}
\usepackage{xcolor}
\usepackage[titles]{tocloft}

\makeatletter
\DeclareRobustCommand{\brsupcite}[1]{\textsuperscript{[\citen{#1}]}}
\makeatother
\let\oldcite\cite
\renewcommand{\cite}[1]{\brsupcite{#1}}

\geometry{top=30mm,bottom=25mm,left=30mm,right=20mm}
\onehalfspacing
\ifnum\ifxetex1\else0\fi=1
    \setmainfont{Times New Roman}[BoldFont={Times New Roman}, ItalicFont={Times New Roman Italic}]
    \setCJKmainfont[BoldFont={SimHei}]{SimSun}
    \setCJKfamilyfont{hei}{SimHei}
\fi
\newcommand{\hei}{\CJKfamily{hei}}
\newcommand{\erhao}{\zihao{2}}
\newcommand{\sanhao}{\zihao{3}}
\newcommand{\sihao}{\zihao{4}}
\newcommand{\xiaosi}{\zihao{-4}}
\newcommand{\xiaowu}{\zihao{-5}}

\titleformat{\section}{\centering\hei\bfseries\sanhao}{\thesection}{1em}{}
\titleformat{\subsection}{\hei\bfseries\sihao}{\thesubsection}{1em}{}
\setcounter{secnumdepth}{3}
\setcounter{tocdepth}{2}
\setlength{\parindent}{2em}
\setlength{\parskip}{0pt}
\pagestyle{fancy}
\fancyhf{}
\fancyfoot[R]{\xiaowu\thepage}
\renewcommand{\headrulewidth}{0pt}
\renewcommand{\cftsecfont}{\hei\bfseries\xiaosi}
\renewcommand{\cftsubsecfont}{\songti\xiaosi}

% 已按 skill 规则将论文封面替换为文献翻译封面
\newcommand{\makecover}{%
      \IfFileExists{/g:/My_code/Agent_atuo_graduation_project_of_GDUT/模板/文献翻译封面.pdf}{\includepdf[pages=-]{/g:/My_code/Agent_atuo_graduation_project_of_GDUT/模板/文献翻译封面.pdf}}{}%
}

\begin{document}
\makecover
\pagenumbering{gobble}
\tableofcontents
\clearpage
\pagenumbering{arabic}
\setcounter{page}{1}

\section{外文文献译文（1）}

\subsection{译文信息}
\begin{itemize}[leftmargin=2em]
    \item 外文题名：Analysis and Comparison of UART, SPI and I2C
    \item 中文译名：UART、SPI 与 I2C 的分析与比较
    \item 作者：未在本次测试中成功提取，需人工复查
    \item 来源/期刊/会议/出版社：未在本次测试中成功提取，需人工复查
    \item 年份：未在本次测试中成功提取，需人工复查
    \item DOI 或 URL：未在本次测试中成功提取，需人工复查
\end{itemize}

\subsection{翻译正文}
本文件仅用于验证 skill 流程，不进行整篇翻译。按 skill 设计，此处应填入从 PDF 首页或前 1--2 页提取的原文段落之中文译文，并按自然段组织，单次处理不超过约 500 字。然而，本次输入文件在读取阶段被识别为受密码保护的 PDF，导致自动提取失败，因此这里只保留流程占位内容以验证模板整合、目录结构与封面替换步骤是否可运行。

可确认的最小信息是题名 ``Analysis and Comparison of UART, SPI and I2C''，其可作为“译文信息”字段的来源。其余作者、来源、年份以及首页正文段落，均需在提供未加密 PDF 或解锁后再继续提取与翻译。

\section{外文文献原文（1）}

\subsection{文献原文}
本节按 skill 规则应保留被翻译的外文原文，并维持原始段落顺序，同时对 LaTeX 特殊字符进行转义。由于本次测试输入 PDF 无法直接读取前 1--2 页文本，故未能自动填入原文，仅记录以下已知信息与失败原因。

Known title: Analysis and Comparison of UART, SPI and I2C.

Extraction status: the source PDF is password-protected, so the first-page or first-two-page text could not be extracted automatically in this test run.

\end{document}
